\section{Overview}
The signal conversion module (SC) converts the analog EEG signal
acquired by the low-level signal processing module to a 12-bit digital
signal representation. The SC module consists of a hardware interface
and AD converter and Windows based PC driver software.

The digital to analog conversion (ADC) of the acquired EEG signal is
not considered a priority module in this system as optimized
commercial ADC's are available. It is however prudent to take
cognizance of signal noise and error sources due to the conversion
process and note their overall contribution to the system's SNR. While
keeping noise and error contribution factors in mind a cost effective
custom made ADC system can be designed and implemented that will have
the minimum impact on signal quality.


The low bandwidth requirement eases the implementation of the ADC
module as high frequency integrated analog to digital converters are
expensive and have strict restrictions on circuit board layout design
and implementation. Additional circuit complexity is introduced by the
approximate DC qualities at the low frequency end of a normal EEG
signal. The low bandwidth of the EEG signal of interest in this project
(0.1~--~35~Hz) makes the use of standard sound cards as ADC's difficult
as most commercial sound cards have a 20~Hz high--pass filter at their
front ends.

Because the design and implementation of a ADC system is fairly
tangential to the topic of active electrode design a commercial 12-bit
SC module is used for data capturing purposes.



\section{Design specification}
As the EEG bandwidth of interest in this project ranges from 0.1~Hz to
35~Hz the AD sampling frequency according to the minimum Nyquist
frequency $f_{nq}$ requirement can be significantly lower than that of
data acquisition systems for high bandwidth applications. This
specification makes it possible to use low--cost and readily available
integrated ADC solutions to implement this module.


Depending on the application a certain amount of signal error is
permissible. Biofeedback applications depend on gross values only and
may be effectively implemented with a low accuracy ADC module. For the
purpose of clinical diagnosis a minimum resolution requirement of 12
bits is usually the case. Increased resolution is not always useful as
amplifiers obtainable with current solid state technology will swamp
the extra bits with noise \cite{pc-eeg}. Non-linear encoding
techniques may be employed to make use of the extra resolution with
the added cost of significant increase in system complexity and
implied expense. A 12--bit DAC is required to insure a high level of
data integrity in the signal display module.

The SC module must interface directly with standard IBM compatible PC
hardware. A standard ECP/ECP parallel port will be used. Making use of
the standard capabilities available in modern PC hardware greatly
reduces the complexity of data acquisition system
implementations. Personal computer software will be used to capture,
store and display the digitized EEG signal.

\section{Implementation}
\subsection{Hardware interface}

The signal conversion module hardware is implemented using the ADC-42
device available from Pico Technology Limited
\footnote{http://www.picotech.com}.

\begin{figure}[htbp]
	\begin{center}
	\includegraphics{adc42.eps} 
	\caption{Picotech ADC-42.}  
	\label{fig:ladc42}
	\end{center}
\end{figure}

Figure~\ref{fig:ladc42} is a photograph of the ADC-42 as found on the
companies Web--site. The ADC-42 is compact single channel 12--bit
analog to digital converter that plugs into a standard PC parallel
port. A single BNC connection is available that can be used with a
standard oscilloscope probe. 

Table~\ref{table:adc42} is a summary of the ADC-42 specification:

\begin{table}
\begin{center}	
	\begin{tabular}[htpb]{|c|l|} \hline
		Resolution & 12~bit \\
		Input voltage range & $\pm$12~V \\
		Input Impedance ($Z_{in}$) & 1~M$\Omega$ \\
		Sample Rate & 15~ksps (kilo samples per second) \\
		Absolute accuracy at 25$^oC$ & $\pm$1\% \\
		\hline		
	\end{tabular}
	\caption{ADC-42 product specification}
	\label{table:adc42}
\end{center}	
\end{table}

From Table~\ref{table:adc42} can be seen that all the requirements of
the design specification are met. The ADC has a resolution of 12 bits
that ensures a good digital representation of the analog EEG
signal. The high input impedance implies that no unnecessary loading
of the filter output stages will occur. The high $Z_{in}$ value
negates the need for a output buffer in the LLSPM. The sample rate of
15~ksps is more than adequate for the (0.1--35~Hz) EEG band of
interest.

The ADC-42 complies with the EEC requirements on generic EMI
emissions. The close proximity of a PC to the low level signal
processing module does however have noise implications. For this
reason commercial oscilloscope cables were used and the LLSP and SA
modules were positioned as far as possible from the PC.

\subsubsection{ADC Safety}
The ADC-42 ground input is connected directly to the PC ground. This is
done to minimize interference. Care must be taken to ensure that the
ground input of the ADC is protected from large voltage
differences. 

\subsection{ADC driver software}

The ADC-42 package is shipped with all the necessary driver software
and tools needed to use the ADC-42 as a standard oscilloscope or data
logger. 

The standard ADC-42 software were extensively used throughout the
testing and measurement phases of the various modules.

\subsection{Conversion noise}

A main concern when using or designing ADC systems are the systems
sampling rate and resolution and the effect of said parameters on
signal quality. The signal conversion process is inherently 'noisy' as
it introduces a conversion error for each digital approximation of the
measured analog signal. The signal error or conversion noise
introduced in this manner can be reduced by increasing the converter's
resolution but not completely eliminated. The acceptable amount of
conversion noise in noted and the system's accuracy predicted based on
the conversion noise figure. It is possible to reduce conversion noise
in a averaging system by assuming that the error value is fairly
constant \cite{roundoff-error}. The ADC-42 driver software includes
facilities to activate this feature at the source level.

The SNR of a quantified signal produced by a sequential approximation
converter, as used in the ADC-42 (PIC16C774), is a function of the
variance of the sampled data:
\begin{equation}
	SNR = 6.02b + 10.79 +10\log_{10}\sigma^2
	\label{eq:cnoise}
\end{equation}
With $b$ the ADC resolution and $\sigma$ the variance of the sampled
data. Equation~\ref{eq:cnoise} assumes no other sources of noise. The
mean of the resulting digital signal is determined over a set of
samples and each sample corrected with the calculated error. Typical
SNR values for the PIC16C774 used in the implementation of the ADC-42
is $\pm\/80~dB$~\footnote{http://www.microchip.com}.

The resolution of the ADC is defined as the smallest voltage
difference that can be detected. The system resolution is a function
of the input voltage range and the converter's resulting word width:
\begin{equation}
	Res = \frac{Range}{ 2^{nbits}- 1}
	\label{eq:resolution}
\end{equation}
With $Res$ the resolution, $Range$ the input voltage range (10~V) and
$nbits$ the word width (12). For the ADC-42 the resolution
$Res$~=~2,4~mV or 0.024\% of full-scale. It is therefore important to
amplify the EEG signal to close to full scale to keep conversion
errors within this acceptable margin.

